\documentclass[a4paper,11pt]{article}
\usepackage[utf8]{inputenc}
\usepackage[T1]{fontenc}
\usepackage[french]{babel}
\usepackage{times}
\usepackage[margin=1.2cm,top=1cm,bottom=1cm]{geometry}
\usepackage{graphicx}
\usepackage{tabularx}
\usepackage{array}
\usepackage{tcolorbox}
\usepackage{enumitem}
\usepackage{xcolor}

\pagestyle{empty}
\setlength{\parindent}{0pt}
\setlength{\parskip}{3pt}

% Define custom dotted rule
\newcommand{\drule}[1]{\makebox[#1]{\dotfill}}

\begin{document}

% =========================================================================
% PAGE 1 : Compétence 01 - Compréhension et communication orales
% =========================================================================

\noindent
\begin{tabularx}{\textwidth}{|p{3cm}|p{3.5cm}|p{3.5cm}|X|}
\hline
Durée: 1Heure & Date: & Niveau: 1ére Am & Année scolaire: 2026/2027 \\
\hline
\end{tabularx}

\vspace{4pt}
\noindent\textbf{Compétence 01: compréhension et communication orales}\\
\textbf{Objectifs:}\\
\_Identifier le théme de la situation de communication.\\
\_Identifier les unités de sens.\\
\_S'exprimer en fonction de la situation de communication.\\
\textbf{Difficultés:}\\
\_L'élève n'arrive pas à identifier le thémé lors de l'écoute d'un suuport audio.\\
\_ L'élève confond entre le théme et les autres informations citées dans le support audio.

\vspace{3pt}
\noindent\textbf{Les élévés qui ont des difficultés :}\\
\_{}\\[-3pt]
\_{}\\[-3pt]
\_{}\\[-3pt]
\_{}\\[-3pt]
\_{}\\[-3pt]
\_{}\\[-3pt]
\_{}\\[-3pt]
\_{}\\[-3pt]
\_{}\\[-3pt]
\_{}

\vspace{-2pt}
\begin{center}
\textbf{Déroulement de l'activité}
\end{center}
\vspace{-4pt}
\_Avez-vous déja visité un zoo avec vos parents ?\\
\_Si oui, dites ce qu'on trouve dedans ?

\vspace{-2pt}
\begin{center}
\textbf{Phase d'identification}
\end{center}
\vspace{-4pt}
\noindent\textbf{texte transcrit :}\\
\vspace{-2pt}
\begin{tcolorbox}[colback=white,colframe=black,arc=0mm,boxrule=0.8pt,top=4pt,bottom=4pt,left=6pt,right=6pt]
La gazelle est un mammifiére sauvage qui vit dans le désert. Elle est herbivore, elle se nourrit
d'herbes, de plantes et de feuilles d'arbes. Son poil est de couleur rousse sur le dos et blanche sur le ventre.
Elle a une queue noire et quatre pattes fines terminées par des sabots.
Sa petite téte porte deux cornes. La gazelle pèse jusqu'a 30 kilos.
\begin{center}
La gazelle est animal trés rapide.
\end{center}
\end{tcolorbox}

\vspace{-2pt}
\noindent\textbf{Premiére écoute du document audio (pré-écoute)}\\
\_Qui parle dans ce vidéo ? \hfill a) un adulte \hfill b) un enfant \hspace{3cm}\null\\
\_ De quoi parle t'il ? \hfill a) du lion \hfill b) de la gazelle \hfill c) du fennec \hspace{1cm}\null\\
\_La gazelle est un animal : \hfill a) sauvage \hfill b) doméstique \hspace{3cm}\null

\vspace{2pt}
\noindent\textbf{Deuxiéme écoute du document audio (écoute)}\\
\_La gazelle est herbivore elle se nourrit : \hfill a) d'herbe \hfill b) de viande \hspace{3cm}\null\\
\_Combien pèse la gazelle?\\
\_Où vit-elle ? \hfill a) dans la savane africaine \hfill b) dans le jardin \hspace{2.5cm}\null

\vspace{2pt}
\noindent\textbf{Troisiéme écoute du document audio (post-écoute)}\\
\_Les pattes de la gazelle sont-il : \hfill a) fines \hfill b) grosses \hspace{3cm}\null\\
\_Sa queue est : \hfill a) noire \hfill b) jaune \hspace{3.5cm}\null\\
\_Comment est décrite la gazelle à la fin de la vidéo ?

\vspace{4pt}
\noindent\textbf{Synthése}\\
Compléte l'énoncé suivant :

\vspace{2pt}
\noindent
\begin{minipage}[c]{0.56\textwidth}
La gazelle est un animal \dotfill ,elle\\[4pt]
est \dotfill , elle pèse environ\\[4pt]
\dotfill ,elle vit dans \dotfill .
\end{minipage}
\hfill
\begin{minipage}[c]{0.40\textwidth}
\centering
\includegraphics[width=\linewidth,height=3.4cm,keepaspectratio]{images/gazelle_p1.png}
\end{minipage}

\newpage
% =========================================================================
% PAGE 2 : Fiche élève (2 colonnes découpables)
% =========================================================================

\newcommand{\gazelleStudentSheet}{%
\begin{minipage}[t][27cm][t]{0.96\linewidth}
\vspace{4pt}
\textbf{Support : La gazelle}

\begin{center}
\textbf{Questionnaire}
\end{center}

\textbf{Premiére écoute :}\\
1/ Qui parle dans ce vidéo?\\
\hspace*{1.5cm} a) un adulte \hfill b) un enfant\hspace*{1.5cm}\null\\[3pt]
2/ De quoi parle t-il ?\\
\hspace*{0.5cm} a) du lion \hfill b) de la gazelle \hfill c) du fennec\hspace*{0.5cm}\null\\[3pt]
3/ La gazelle est un animal :\\
\hspace*{1.5cm} a) sauvage \hfill b) doméstique\hspace*{1.5cm}\null

\vspace{4pt}
\textbf{Deuxiéme écoute :}\\
1/ La gazelle est un herbivore elle se nourrit :\\
\hspace*{1.5cm} a) d'herbe \hfill b) de la viande\hspace*{1.5cm}\null\\[3pt]
2/ Où vit la gazelle ?\\
\hspace*{1.5cm} a) dans la savane \hfill b) dans le jardin\hspace*{1.5cm}\null\\[3pt]
3/ Combien pèse la gazelle ?

\vspace{4pt}
\textbf{Troisiéme écoute :}\\
1/ Les pattes de gazelle sont -ils :\\
\hspace*{1.5cm} a) fines \hfill b) grosses\hspace*{1.5cm}\null\\[3pt]
2/ Sa queue est : \hspace{1cm} a) noire \hfill b) jaune\hspace*{1.5cm}\null\\[3pt]
3/ Comment est décrite la gazelle à la fin de la vidéo?\\[3pt]
\dotfill

\vspace{6pt}
\noindent\makebox[\linewidth]{\dotfill}

\vspace{4pt}
\textbf{Synthése}\\
Compléte l'énoncé suivant :\\
La gazelle est un animal \dotfill , elle\\[4pt]
est \dotfill , elle se nourrit de\\[4pt]
\dotfill , elle vit dans la\\[4pt]
\dotfill .

\vspace{10pt}
\begin{center}
\includegraphics[width=\linewidth,height=8cm,keepaspectratio]{images/gazelle_schema.png}
\end{center}
\end{minipage}%
}

\noindent
\begin{tabular}{|p{0.485\textwidth}|p{0.485\textwidth}|}
\hline
\gazelleStudentSheet & \gazelleStudentSheet \\
\hline
\end{tabular}

\newpage
% =========================================================================
% PAGE 3 : Compétence 02 - Lecture
% =========================================================================

\noindent
\begin{tabularx}{\textwidth}{|p{4cm}|p{4.5cm}|X|}
\hline
Durée: 1Heure & Date: & Niveau: 1ére Am \\
\hline
\end{tabularx}

\vspace{4pt}
\noindent\textbf{Compétence 02: Lécture}\\
\textbf{Critère :}\\
\_Activer la correspondance graphie/phonie.\\
\_Réaliser une lécture fluide.\\
\_Réspecter l'intonnation.\\
\textbf{Objectifs :} L'éléve sera capable de dévélopper la compétence de la lécture.\\
\_ Lire de maniére courante et expressive\\
\_ Distinguer entre les sons: /g/ /3/, /s/ /z/. /k/ /s/\\
\textbf{Description des obstacles :}\\
\_ L'éléve ne sait pas lire la lettre qui correspond au phonéme.\\
\_ L'élève n'arrive pas à lire correctement un énoncé.\\
\_ La majorité des élèves ont du mal à réspecter l'intonation des types de phrases.

\vspace{3pt}
\noindent\textbf{Les éléves qui ont des difficultés :}\\
\_{}\\[-3pt]
\_{}\\[-3pt]
\_{}\\[-3pt]
\_{}\\[-3pt]
\_{}\\[-3pt]
\_{}\\[-3pt]
\_{}\\[-3pt]
\_{}\\[-3pt]
\_{}\\[-3pt]
\_{}

\vspace{3pt}
\noindent\textbf{Modalité / Groupes:} Travail en groupes.\\
\textbf{Matériel / Support:} Tableau, cahier d'exercice

\vspace{4pt}
\begin{tcolorbox}[colback=white,colframe=black,arc=0mm,boxrule=0.8pt,halign=center,top=3pt,bottom=3pt]
\textbf{Mise en oeuvre de la remédiation}
\end{tcolorbox}

\vspace{-2pt}
\noindent\textbf{Présentation:}

\vspace{2pt}
\noindent
\begin{tabularx}{\textwidth}{|X|X|}
\hline
\textbf{Role de l'enseignante} & \textbf{Role des apprenants} \\
\hline
\_Présente les activités. & \_ Lisent les activités. \\
\_Explique les consignes. & \_Ecoutent en posent des questions. \\
\_Oriente et guide les apprenants. & \_ S'entraident et échangent entre eux. \\
\hline
\end{tabularx}

\vspace{4pt}
\noindent\textbf{Moment de Mise en confiance :}\\
\textbf{Devinette qui suis-je?}\\
Je travaille à l'hopital, je soigne les malades, j'utilise le stéthoscope, qui suis-je?

\vspace{4pt}
\noindent\textbf{\textit{\underline{Activités proposées}}}

\vspace{2pt}
\noindent\textbf{Activité 01:} Lis attentivement le passage suivant:\\
Chaque jours, Amira va à l'école à 8 Heures du matin. Dans la classe, elle retrouve ses camarades et écoute sa maitresse. elle apprend à lire et à écrire, elle aime dessiner les fleurs et les animaux.

\vspace{3pt}
\noindent\textbf{Activité 02:} Classe les mots suivants dans le tableau selon la pronociation de la lettre ``C''\\
Carotte-cerise-coq-bicyclette-cube-citron

\vspace{2pt}
\noindent
\begin{tabularx}{\textwidth}{|X|X|}
\hline
\textbf{/k/} & \textbf{/S/} \\
\hline
\rule{0pt}{1.8em} & \\
\hline
\end{tabularx}

\vspace{3pt}
\noindent\textbf{Activité 03:} Prépare deux enveloppes: un pour le son /g/ et l'autre pour le son /3/\\
Mettez chaque mots dans un enveloppes adéquat: Singe-gaspiller-guide-garçon-plage-gomme.

\vspace{2pt}
\noindent\textbf{Activité 03:} Lis les phrases suivantes correctement et respecte l'intonation\\
\_ Les enfants visitent le zoo.\\
\_ aimez-vous les animaux?\\
\_ quel beau métier!

\newpage
% =========================================================================
% PAGE 4 : Compétence 03 - Compréhension de l'écrit (Partie 1)
% =========================================================================

\noindent
\begin{tabularx}{\textwidth}{|p{3cm}|p{3.5cm}|p{3.5cm}|X|}
\hline
Durée: 1Heure & Date: & Niveau: 1ére AM & Année scolaire: 2026/2027 \\
\hline
\end{tabularx}

\vspace{4pt}
\noindent\textbf{Compétence 03: Compréhension de l'écrit}\\
\textbf{Critére 1 et 2 :} Identifier le théme général et le champs léxical relatif au texte\\
\textbf{Difficultés :}\\
\_L'éléve trouve des difficultés à repérer le théme\\
\_ l'éléve peine à repérer les mots en relation avec le théme

\vspace{3pt}
\noindent\textbf{Les éléves qui ont des difficultés :}\\
\_{}\\[-3pt]
\_{}\\[-3pt]
\_{}\\[-3pt]
\_{}\\[-3pt]
\_{}\\[-3pt]
\_{}\\[-3pt]
\_{}\\[-3pt]
\_{}\\[-3pt]
\_{}\\[-3pt]
\_{}\\[-3pt]
\_{}

\vspace{3pt}
\noindent\textbf{Modalité de travail :} En groupe puis individuel

\vspace{4pt}
\noindent\textbf{I)- Mise en situation :}\\
* Proposer aux élève de citer quelques noms de métiers\\
* Que porte un medecin ? Blouse blance (tenuque médicale) (Matériels médicales)

\vspace{4pt}
\noindent\textbf{II)- Phase d'identification :}\\
* Proposer aux éleves trois textes (sans titres):

\vspace{2pt}
\begin{tcolorbox}[colback=white,colframe=black,arc=0mm,boxrule=0.8pt,top=3pt,bottom=3pt,left=5pt,right=5pt]
\textbf{Texte 1 :}\\
Le pompier intervient lors les incendies pour éteindre le feu et sauver les personnes en danger.Il utilise des tuyaux, des lances à eau et porte une tenue spéciale pour se protéger. Le pompier aide aussi lors des accidents et des innondations. Le métier de pompier est trés noble car il permet de sauver des vies.
\end{tcolorbox}

\vspace{-3pt}
\begin{tcolorbox}[colback=white,colframe=black,arc=0mm,boxrule=0.8pt,top=3pt,bottom=3pt,left=5pt,right=5pt]
\textbf{Texte 2 :}\\
Le lapin est un animal doméstique. Il est doux et appréciè des enfants. Il se nourrit de carottes, de légumes et d'herbe. Il vit dans une cage ou dans un clapier et peut etre élèvé à la ferme ou dans la maison.
\end{tcolorbox}

\vspace{-3pt}
\begin{tcolorbox}[colback=white,colframe=black,arc=0mm,boxrule=0.8pt,top=3pt,bottom=3pt,left=5pt,right=5pt]
\textbf{Texte3 :}\\
L'environnement est un lieu qui nous entoure. Pour vivre sainement, nous devons protéger ce milieu de tout ce qui peut le poluer. Pour cela, la participation de tous les citoyens est nécessaire.
\end{tcolorbox}

\vspace{2pt}
\noindent\_ Les élèves attribuent les titres et les images aux textes qui corresponden

\vspace{2pt}
\noindent
\begin{tabularx}{\textwidth}{XXX}
* Le lapin & * Le pompier & * L'environnement \\
* Doméstique & * Polluer & * Sauver \\
\end{tabularx}

\vspace{4pt}
\noindent
\begin{minipage}{0.32\textwidth}
\centering
\includegraphics[height=3.4cm,keepaspectratio]{images/pompier.png}
\end{minipage}
\hfill
\begin{minipage}{0.32\textwidth}
\centering
\includegraphics[height=3.4cm,keepaspectratio]{images/lapin.png}
\end{minipage}
\hfill
\begin{minipage}{0.32\textwidth}
\centering
\includegraphics[height=3.4cm,keepaspectratio]{images/environnement.png}
\end{minipage}

\newpage
% =========================================================================
% PAGE 5 : Compétence 03 - Compréhension de l'écrit (Partie 2)
% =========================================================================

\noindent\textbf{III)- Phase d'application :}\\
\textbf{Activité1 :}\\
Proposer un texte sans titre ni image puis demander aux élèves de sélectionner le titre qui convient le mieux au texte parmi la liste composée de trois propositions :

\vspace{4pt}
\begin{tcolorbox}[colback=white,colframe=black,arc=0mm,boxrule=0.8pt,top=4pt,bottom=4pt,left=6pt,right=6pt]
\textbf{Texte :}\\
La girafe est animal sauvage. Elle est le plus grand animal terrestre. Elle se nourrit de feuilles et de fruits qu'elle cueille en haut des arbres grace à son long cou. Elle vit dans la savane et se déplace en troupeaux.
\end{tcolorbox}

\vspace{4pt}
\noindent\textbf{* Parmi ces titres lequel convient le mieux au texte :}\\
\_ Le titre du texte est :\\
\_ Le lapin\\
\_ La girafe\\
\_ Le tremblement de terre

\vspace{6pt}
\noindent\textbf{*Releve du texte le champ léxicale `` Animal''}\\[4pt]
\noindent\makebox[\linewidth]{\dotfill}

\vspace{6pt}
\noindent\textbf{IV)- Phase de ré investissement:}\\
\_ Compléte ce texte à l'aide de ces mots : animaux \_ grimaces \_ girafes \_ pain\\[4pt]
\noindent\makebox[\linewidth]{\dotfill}

\vspace{4pt}
\noindent J'ai visité le Zoo avec ma famille. J'ai vu beaucoup d'\dotfill{} Les singes m'ont beaucoup amusé: ils faisaient des \dotfill{} J'ai donné à manger à la \dotfill{} `` IL est interdit de donner du \dotfill{} aux animaux `` dit le gardien.

\vspace{10pt}
\noindent
\begin{minipage}{0.48\textwidth}
\centering
\includegraphics[width=\linewidth,height=10cm,keepaspectratio]{images/girafe_arbre.png}
\end{minipage}
\hfill
\begin{minipage}{0.48\textwidth}
\centering
\includegraphics[width=\linewidth,height=10cm,keepaspectratio]{images/girafe_savane.png}
\end{minipage}

\newpage
% =========================================================================
% PAGE 6 : Compétence 04 - Production Écrite (Fiche)
% =========================================================================

\noindent
\begin{tabularx}{\textwidth}{|p{3cm}|p{3.5cm}|p{3.5cm}|X|}
\hline
Durée: 1Heure & Date: & Niveau:ére 1AM & Année scolaire: 2026/2027 \\
\hline
\end{tabularx}

\vspace{4pt}
\noindent\textbf{Compétence 04: Production Écrite}\\
\textbf{Objectifs:}\\
\_ Pertinence du texte produit.\\
\_ Cohérence du texte produit.\\
\_ Correction de la langue.\\
\_ Lisibilité de l'écrit (qualité de la graphie).

\vspace{3pt}
\noindent\textbf{Les éléves qui ont des difficultés :}\\
\_{}\\[-3pt]
\_{}\\[-3pt]
\_{}\\[-3pt]
\_{}\\[-3pt]
\_{}\\[-3pt]
\_{}\\[-3pt]
\_{}\\[-3pt]
\_{}\\[-3pt]
\_{}\\[-3pt]
\_{}

\vspace{3pt}
\noindent\textbf{Sujet :}\\
À partir de la fiche d'identité du lion, écris un court texte de quatre à cinq phrases dans lequel tu présente `` le lion `` .\\
\_ Donner un titre.\\
\_ Utiliser des phrases courtes.\\
\_ Ponctue ton texte.

\vspace{4pt}
\begin{tcolorbox}[colback=white,colframe=black,arc=0mm,boxrule=0.8pt,top=4pt,bottom=4pt,left=6pt,right=6pt]
\begin{minipage}[c]{0.58\textwidth}
\textbf{Fiche d'identité du lion}\\[3pt]
\textbf{Nom :} Lion\\
\textbf{La femelle :} Lionne\\
\textbf{Le petit :} Linceau\\
\textbf{Il vit :} Dans la savane d'afrique\\
\textbf{Ce qu'il mange :} De la viande `` c'est un carnivore''\\
\textbf{Poids :} jusqu'au 200 kg\\[4pt]
\textbf{Son corps :} Grand et fort
\end{minipage}
\hfill
\begin{minipage}[c]{0.38\textwidth}
\centering
\includegraphics[width=\linewidth,height=3.8cm,keepaspectratio]{images/lion_corps.png}
\end{minipage}
\end{tcolorbox}

\vspace{8pt}
\noindent
\begin{minipage}[t]{0.64\textwidth}
\textbf{Rédaction :}\\[6pt]
\centerline{\makebox[0.55\linewidth]{\dotfill}}\\[14pt]
\noindent\makebox[\linewidth]{\dotfill}\\[14pt]
\noindent\makebox[\linewidth]{\dotfill}\\[14pt]
\noindent\makebox[\linewidth]{\dotfill}\\[14pt]
\noindent\makebox[\linewidth]{\dotfill}\\[14pt]
\noindent\makebox[\linewidth]{\dotfill}
\end{minipage}
\hfill
\begin{minipage}[t]{0.33\textwidth}
\centering
\includegraphics[width=\linewidth,height=6.5cm,keepaspectratio]{images/lion_famille.png}
\end{minipage}

\newpage
% =========================================================================
% PAGE 7 : Fiche élève - 6 étiquettes de Production Écrite (Lion)
% =========================================================================

\newcommand{\lionCard}{%
\begin{minipage}[t][8.3cm][t]{0.96\linewidth}
\vspace{2pt}
\begin{minipage}[t]{0.58\linewidth}
\textbf{Production écrite}\\[2pt]
\textbf{Sujet :}\\
{\fontsize{8.5pt}{10.5pt}\selectfont
À partir de la fiche d'identité du lion, écris un court texte de quatre à cinq phrases dans lequel tu présente `` le lion `` .\\[2pt]
\_ Donner un titre.\\
\_ Utiliser des phrases courtes.\\
\_ Ponctue ton texte.\\[2pt]
\textbf{Fiche d'identité du lion}\\
Nom : Lion\\
La femelle : Lionne\\
Le petit : Linceau\\
Il vit : Dans la savane d'afrique\\
Ce qu'il mange : De la viande `` c'est un carnivore''\\
Poids : jusqu'au 200 kg\\
Son corps : Grand et fort
}
\end{minipage}\hfill
\begin{minipage}[t]{0.38\linewidth}
\vspace{4pt}
\centering
\includegraphics[width=\linewidth,height=6.8cm,keepaspectratio]{images/lion_marche.png}
\end{minipage}
\end{minipage}%
}

\noindent
\begin{tabular}{|p{0.485\textwidth}|p{0.485\textwidth}|}
\hline
\lionCard & \lionCard \\
\hline
\lionCard & \lionCard \\
\hline
\lionCard & \lionCard \\
\hline
\end{tabular}

\newpage
% =========================================================================
% PAGE 8 : Fiche élève - Vignettes schémas de la Gazelle (15 étiquettes)
% =========================================================================

\newcommand{\gazelleSticker}{%
\begin{minipage}[c][4.9cm][c]{\linewidth}
\centering
\includegraphics[width=0.94\linewidth,height=4.6cm,keepaspectratio]{images/gazelle_schema.png}
\end{minipage}%
}

\noindent
\begin{tabular}{|p{0.315\textwidth}|p{0.315\textwidth}|p{0.315\textwidth}|}
\hline
\gazelleSticker & \gazelleSticker & \gazelleSticker \\
\hline
\gazelleSticker & \gazelleSticker & \gazelleSticker \\
\hline
\gazelleSticker & \gazelleSticker & \gazelleSticker \\
\hline
\gazelleSticker & \gazelleSticker & \gazelleSticker \\
\hline
\gazelleSticker & \gazelleSticker & \gazelleSticker \\
\hline
\end{tabular}

\end{document}
